\section{別途与えられたMF同値の残余圏への移送}\label{ar:sec:application}

本節では、非可換MF圏と可換MF圏の比較そのものを構成しない。
その比較を明示的な入力にして、前節までの結果の適用を述べる。

\begin{corollary}[互換なMF比較が与えられた場合]\label{ar:cor:conditional}
$d\ge5$とし、各$N$で第\ref{am:sec:comparison}節の入力$(A_N,\Omega_N)$を取る。
滑らかな特殊ファイバーを持つ形式三次式$W_h\in\mathbb C[[h]][z_1,\ldots,z_d]_3$と、
内部shift・三次周期・Artin商への底変更を保つ同値
\[
 \MF^{\ZZ}(A_N,\Omega_N)\simeq
 \MF^{\ZZ}(R_N[z_1,\ldots,z_d],W_h\bmod h^{N+1})
\]
が各$N$について与えられ、それらが互換であると仮定する。
$B_N=A_N/(\Omega_N)$、$Y_N=V(W_h\bmod h^{N+1})$とすると、互換なpair同値
\[
 (\mathcal K_{B_N}^{\rm lin},\rho)\simeq
 (\Ku_{\rm perf}(Y_N/R_N),\rho_{Y_N})
\]
を得る。ここで右辺は$\operatorname{Perf}(Y_N)$から
$\mathcal O,\ldots,\mathcal O(d-4)$を除いた右直交である。
\end{corollary}
\begin{proof}
定理\ref{am:thm:comparison}を両側に適用し、その間に仮定したMF同値を挟む。
可換側では、線束$\mathcal O(-j)$（$0\le j\le d-1$）の直和が
$D_{\rm QCoh}(Y_N)$をcompactに生成する\cite[Tag~0BQT]{AMStacks}。
これは射影空間の相対対角分解と、閉埋込みのpushforwardの保守性・随伴による。
qcqsスキームではcompact対象とperfect対象が一致する\cite[Tag~09M1]{AMStacks}ので、
全線束の厚い閉包は$\operatorname{Perf}(Y_N)$である。
各比較のrotationと底変更の両立性が、所要の互換性を与える。
\end{proof}

この系は、$W_h$やMF同値の存在証明を代用しない。
非零の複素パラメータへの代数化、完備底での生成性、
全coherent圏との同定についても別の議論が必要である。

\begin{example}[絶対特異圏を代用できない理由]\label{cg:ex:artin-singularity}
$D=\mathbb C[\varepsilon]/(\varepsilon^2)$、$A=D[x]$、$W=x$とする。
有限自由$\mathbb Z/2$-graded MF行列$Q$について$Q^2=xI$を微分すると
$Q(\partial_xQ)+(\partial_xQ)Q=I$となるので、このMF圏は可縮である。
他方、$A/(W)=D$の剰余体は$\varepsilon$による無限周期自由分解を持ち、
その自己Extは任意の非負次数で非零である。
従って剰余体はperfectでなく、絶対特異圏$D_{\rm sg}(D)$は零でない。
周囲でperfectという条件を外した比較は、この例で既に成立しない。
\end{example}
